
\chapter{Diagramas de Secuencia}\label{ch:diagramassecuencia}

Los diagramas de secuencia modelan la interacción entre objetos a lo largo del tiempo para un flujo específico del sistema. A diferencia del diagrama de robustez (Capítulo~\ref{ch:diagramarobustez}), que identifica las clases de software participantes, los diagramas de secuencia detallan el \emph{orden exacto} de los mensajes, las respuestas y los fragmentos combinados (loops, alternativas, opciones) que ocurren durante la ejecución de un caso de uso.

Se presentan cinco flujos críticos de FisioManager: el registro de feedback post-rutina, la asignación de rutinas a pacientes, la gestión de citas desde el rol de Administrador, la mensajería en tiempo real y el auto-agendamiento de cita por el paciente.

\section{Flujo 1: Registro de Feedback}\label{sec:sec-feedback}

Este diagrama modela la secuencia completa desde que el paciente captura su feedback tras una rutina; cuando graba audio, la transcripción y el resumen con IA se generan en ese momento y se guardan junto con el feedback.

\subsection{Participantes}

\begin{itemize}
  \item \actor{Paciente}: actor que completa la rutina e inicia el registro de feedback.
  \item \codeinline{UIRegistroFeedback}: interfaz (boundary) con selector de estado emocional, slider de dolor, mapa corporal y grabador de audio.
  \item \codeinline{feedback.controller}: adapta la petición/respuesta y coordina la subida de audio.
  \item \codeinline{ai.service}: ejecuta la transcripción con Whisper y el resumen clínico con Ollama.
  \item \codeinline{storage}: guarda el archivo de audio (disco local o MinIO/S3).
  \item \codeinline{feedback.service}: valida y persiste el feedback y los puntos de dolor; dispara la alerta de dolor alto.
  \item \codeinline{notification.service}: notifica al terapeuta (in-app y correo) cuando hay dolor alto.
\end{itemize}

\subsection{Descripción del flujo}

\begin{enumerate}
  \item El paciente selecciona su estado emocional, el nivel de dolor (1--10) y marca el mapa corporal de dolor.
  \item Opcionalmente graba una nota de voz. Al subirla (\codeinline{POST /feedback/audio}), se verifica el consentimiento de audio, el \codeinline{storage} guarda el archivo y \codeinline{ai.service} transcribe con Whisper y resume con Ollama, devolviendo \codeinline{audioUrl}, \codeinline{transcripcion} y \codeinline{resumenIA}.
  \item El paciente confirma y envía (\codeinline{POST /feedback}) con los datos numéricos, los puntos de dolor y, si los hubo, la transcripción y el resumen.
  \item El \codeinline{feedback.service} valida (nivel de dolor en rango, estado emocional válido) y persiste el \codeinline{Feedback} con su transcripción y resumen en línea.
  \item Si se marcaron zonas de dolor, se persisten los \codeinline{PuntoDolor} asociados.
  \item Si el nivel de dolor es alto ($\geq 7$), el \codeinline{notification.service} genera la alerta para el terapeuta.
  \item Se confirma al paciente que el feedback fue registrado.
\end{enumerate}

\begin{figure}[H]
  \centering
  \includegraphics[width=\textwidth,height=0.88\textheight,keepaspectratio]{FisioManager_Secuencia_RegistroFeedback.png}
  \caption{Diagrama de Secuencia -- Flujo 1: Registro de Feedback con transcripción y resumen de IA.}
  \label{fig:secuencia-feedback}
\end{figure}

\subsection{Aspectos destacados}

\begin{itemize}
  \item \textbf{IA en la subida del audio:} la transcripción (Whisper) y el resumen (Ollama) se obtienen al subir el audio y se almacenan dentro del propio \codeinline{Feedback}; no existe una entidad ni un procesamiento asíncrono separados.
  \item \textbf{Fragmento \codeinline{opt} para audio:} la grabación es opcional; sin audio el feedback se registra solo con los datos numéricos y el mapa de dolor.
  \item \textbf{Consentimiento informado:} no se procesa audio sin un consentimiento previo registrado con marca de tiempo.
  \item \textbf{Alerta de dolor alto:} el \codeinline{feedback.service} dispara, vía \codeinline{notification.service}, una alerta in-app y por correo al terapeuta cuando el dolor es elevado.
\end{itemize}

\section{Flujo 2: Asignación de Rutina}\label{sec:sec-asignacion}

Este diagrama modela la secuencia completa desde que el terapeuta construye una rutina en el builder hasta que cada paciente seleccionado recibe la notificación de su rutina asignada.

\subsection{Participantes}

\begin{itemize}
  \item \actor{Terapeuta}: actor que crea y asigna rutinas de ejercicio.
  \item \codeinline{UIRutinaBuilder}: interfaz de usuario (boundary) que permite construir rutinas seleccionando ejercicios de la biblioteca, configurando parámetros y asignando a pacientes.
  \item \codeinline{GestorPlanTratamiento}: controlador (control) que orquesta la creación de rutinas y delega la asignación.
  \item \codeinline{AsignacionRutina}: servicio que ejecuta la clonación de la rutina template para cada paciente seleccionado.
  \item \codeinline{Notificación al Paciente}: servicio encargado de generar y enviar las notificaciones a los pacientes.
\end{itemize}

\subsection{Descripción del flujo}

El flujo se divide en dos fases claramente diferenciadas:

\subsubsection{Fase 1: Construcción de la rutina}

\begin{enumerate}
  \item El terapeuta accede al constructor de rutinas.
  \item El sistema carga las plantillas de ejercicio disponibles desde la biblioteca.
  \item El terapeuta selecciona ejercicios aplicando filtros por parte del cuerpo (rodilla, espalda, hombro, cuello, brazo, cadera, tobillo) y por tipo (Físico, Respiración).
  \item Configura los parámetros de cada actividad: duración en minutos, número de repeticiones y tiempo de descanso en segundos.
  \item Confirma la creación de la rutina.
  \item El \codeinline{GestorPlanTratamiento} valida los datos y persiste la rutina como template (\codeinline{patientId = null}) con todas sus actividades.
\end{enumerate}

\subsubsection{Fase 2: Asignación a pacientes}

\begin{enumerate}
  \item El terapeuta selecciona la rutina desde la biblioteca y elige la opción ``Asignar a Pacientes''.
  \item Selecciona uno o más pacientes de su lista.
  \item El \codeinline{GestorPlanTratamiento} delega en \codeinline{AsignacionRutina}.
  \item Para cada paciente seleccionado, \codeinline{AsignacionRutina}:
  \begin{itemize}
    \item Obtiene la rutina template con todas sus actividades.
    \item Clona la rutina asignando el \codeinline{patientId} correspondiente.
    \item Copia todas las actividades con sus parámetros configurados.
    \item Envía una notificación al paciente a través del servicio de notificaciones.
  \end{itemize}
  \item Se confirma al terapeuta que la rutina fue asignada a $N$ pacientes.
\end{enumerate}

\begin{figure}[H]
  \centering
  \includegraphics[width=\textwidth,height=0.88\textheight,keepaspectratio]{FisioManager_Secuencia_AsignacionRutina.png}
  \caption{Diagrama de Secuencia -- Flujo 2: Asignación de Rutina a pacientes.}
  \label{fig:secuencia-asignacion}
\end{figure}

\subsection{Aspectos destacados}

\begin{itemize}
  \item \textbf{Patrón Template--Clone:} la rutina se almacena primero como plantilla reutilizable y luego se clona por cada paciente. Esto permite asignar la misma rutina a múltiples pacientes sin duplicar la configuración manual.
  \item \textbf{Fragmento \codeinline{loop}:} la clonación y notificación se ejecutan dentro de un bucle que itera por cada paciente seleccionado, reflejando la operación por lotes (\textit{batch}).
  \item \textbf{Separación de fases:} el diagrama distingue claramente la fase de construcción (creativa) de la fase de asignación (operativa), lo que facilita la trazabilidad y la posibilidad de reutilizar templates existentes sin reconstruirlos.
  \item \textbf{Notificación integrada:} cada asignación dispara automáticamente una notificación al paciente, garantizando que esté informado de su nueva rutina.
\end{itemize}

\section{Flujo 3: Gestión de Citas (Administrador)}\label{sec:sec-citas}

Este diagrama modela la secuencia completa desde que el Administrador consulta la agenda del día, crea una nueva cita validando la disponibilidad del terapeuta, hasta el envío de recordatorios al paciente. Refleja el rol simplificado de Administrador como único responsable operativo de la clínica.

\subsection{Participantes}

\begin{itemize}
  \item \actor{Administrador}: actor único que gestiona la agenda de toda la clínica; concentra las funciones operativas y técnicas.
  \item \codeinline{UIAgenda}: interfaz de usuario (boundary) que presenta la vista de agenda con filtros por fecha, terapeuta y estado, así como el formulario de creación de citas.
  \item \codeinline{GestorCitas}: controlador (control) que orquesta la validación, creación y consulta de citas.
  \item \codeinline{ValidadorDisponibilidad}: controlador (control) que verifica que el horario solicitado esté dentro del horario del terapeuta y que no existan colisiones con citas ya agendadas.
  \item \codeinline{Cita}: entidad persistente que representa una cita agendada con estado (SCHEDULED, COMPLETED, CANCELLED).
  \item \codeinline{notification.service}: servicio que genera la \codeinline{Notificacion} in-app y, mediante \codeinline{email.service}, envía el recordatorio por correo al paciente.
\end{itemize}

\subsection{Descripción del flujo}

El flujo se divide en tres fases:

\subsubsection{Fase 1: Consulta de agenda}

\begin{enumerate}
  \item El Administrador accede a la agenda del día.
  \item El \codeinline{GestorCitas} recupera las citas filtradas por fecha, terapeuta y/o estado.
  \item La \codeinline{UIAgenda} muestra la lista de citas con estadísticas del día (total, programadas, completadas, canceladas).
\end{enumerate}

\subsubsection{Fase 2: Creación de nueva cita}

\begin{enumerate}
  \item El Administrador completa el formulario de nueva cita (paciente, terapeuta, fecha/hora, motivo, modalidad).
  \item El \codeinline{GestorCitas} valida los campos requeridos y delega en el \codeinline{ValidadorDisponibilidad}.
  \item El \codeinline{ValidadorDisponibilidad} verifica dos condiciones:
  \begin{itemize}
    \item Que la fecha/hora esté dentro del horario registrado del terapeuta.
    \item Que no existan colisiones con citas previamente agendadas en estado SCHEDULED.
  \end{itemize}
  \item Si la disponibilidad es válida:
  \begin{itemize}
    \item Se persiste la cita con estado SCHEDULED (con \codeinline{grupoRecurrencia} si se trata de una serie).
    \item Queda registrada para el recordatorio automático de 24~horas, que un planificador en \codeinline{notification.service} dispara más tarde.
    \item Se confirma al Administrador la creación exitosa.
  \end{itemize}
  \item Si la disponibilidad no es válida, se informa al Administrador el motivo (colisión o fuera de horario) para que seleccione otro horario.
\end{enumerate}

\subsubsection{Fase 3: Envío manual de recordatorio}

\begin{enumerate}
  \item El Administrador selecciona una cita en estado SCHEDULED y solicita enviar el recordatorio.
  \item El \codeinline{GestorCitas} (\codeinline{sendReminder}) recupera los datos de la cita, paciente y terapeuta.
  \item Delega en \codeinline{email.service}, que arma el contenido y lo envía por correo electrónico.
  \item Se confirma al Administrador que el recordatorio fue enviado al paciente.
\end{enumerate}

\begin{figure}[H]
  \centering
  \includegraphics[width=\textwidth,height=0.88\textheight,keepaspectratio]{FisioManager_Secuencia_GestionCitas.png}
  \caption{Diagrama de Secuencia -- Flujo 3: Gestión de Citas por el Administrador.}
  \label{fig:secuencia-citas}
\end{figure}

\subsection{Aspectos destacados}

\begin{itemize}
  \item \textbf{Rol de Admin simplificado:} el Administrador es el único responsable de la gestión de la agenda de toda la clínica, tal como se define en el modelo de actores (Capítulo~2). No se contemplan múltiples administradores.
  \item \textbf{Validación de disponibilidad:} antes de crear cualquier cita, el \codeinline{ValidadorDisponibilidad} verifica tanto el horario del terapeuta como las colisiones con citas existentes, evitando sobreposiciones.
  \item \textbf{Fragmento \codeinline{alt}:} modela la bifurcación entre el flujo exitoso (horario disponible) y el flujo alternativo (colisión o fuera de horario), garantizando que el Administrador reciba retroalimentación clara.
  \item \textbf{Recordatorio automático y manual:} un planificador en \codeinline{notification.service} envía el recordatorio automático 24~h antes (y marca la cita como recordada); adicionalmente, el Administrador puede disparar el envío manual por correo en cualquier momento.
  \item \textbf{Trazabilidad con el diagrama de robustez:} los participantes corresponden directamente a las clases de software identificadas en el diagrama de robustez (Capítulo~\ref{ch:diagramarobustez}): \codeinline{UIAgenda}~(boundary), \codeinline{GestorCitas} y \codeinline{ValidadorDisponibilidad}~(control), \codeinline{Cita}~(entity) y \codeinline{GestorNotificaciones}~(control).
\end{itemize}

\section{Secuencia: Mensajería en tiempo real}

Este diagrama modela el envío de un mensaje del paciente a su terapeuta. La API valida que el emisor sea participante de la conversación, persiste el mensaje y lo entrega al receptor en tiempo real mediante Socket.IO, generando además una notificación in-app. Con el \textit{adapter} de Redis los eventos se propagan entre varias instancias de la API.

\begin{figure}[H]
  \centering
  \includegraphics[width=\textwidth,height=0.88\textheight,keepaspectratio]{FisioManager_Secuencia_Mensajeria.png}
  \caption{Secuencia de mensajería en tiempo real (paciente $\rightarrow$ terapeuta).}
  \label{fig:seq-mensajeria}
\end{figure}

\section{Secuencia: Auto-agendamiento de cita}

Este diagrama modela cómo el paciente consulta las franjas libres de su terapeuta asignado y agenda por sí mismo una cita. La API calcula los horarios disponibles (disponibilidad menos colisiones) y, al confirmar, valida que el terapeuta sea el asignado, que la fecha esté dentro del horario y que no exista conflicto, notificando al terapeuta.

\begin{figure}[H]
  \centering
  \includegraphics[width=\textwidth,height=0.88\textheight,keepaspectratio]{FisioManager_Secuencia_AutoAgenda.png}
  \caption{Secuencia de auto-agendamiento de cita por el paciente.}
  \label{fig:seq-autoagenda}
\end{figure}

\section{Código PlantUML de los diagramas}

Los códigos fuente se entregan como \codeinline{10\_secuencia\_asignacion\_rutina.puml}, \codeinline{11\_secuencia\_registro\_feedback.puml}, \codeinline{13\_secuencia\_gestion\_citas.puml}, \codeinline{14\_secuencia\_mensajeria.puml} y \codeinline{15\_secuencia\_autoagenda.puml} en el directorio \codeinline{puml/}. Se regeneran con:

\begin{lstlisting}[language=bash]
plantuml -tpng 10_secuencia_asignacion_rutina.puml
plantuml -tpng 11_secuencia_registro_feedback.puml
plantuml -tpng 13_secuencia_gestion_citas.puml
plantuml -tpng 14_secuencia_mensajeria.puml
plantuml -tpng 15_secuencia_autoagenda.puml
\end{lstlisting}

El código fuente completo se incluye en el anexo de código fuente PlantUML del Modelo de Análisis.
